\documentclass[conference]{IEEEtran}

\usepackage{graphicx}
\usepackage{booktabs}
\usepackage{url}
\usepackage[hidelinks]{hyperref}

\begin{document}

\title{Predicting 30-Day Hospital Readmission Among Patients with Diabetes}

\author{
\IEEEauthorblockN{Siddique Chand, Aubin Mugisha, and Cameron Letendre}
\IEEEauthorblockA{
University of Southern Maine, Portland, ME, USA
}
}

\maketitle


% =========================================================
\begin{abstract}

% Briefly summarize:
% - What problem did we study?
% - What dataset did we use?
% - What models did we test?
% - What was the main result?


\end{abstract}


% =========================================================
\section{Introduction}

% Explain:
% - What is 30-day hospital readmission?
% - Why is it important?
% - Why are we focusing on patients with diabetes?
% - What is the goal of this project?
% - What are we trying to predict?

[Brief introduction to the problem.]

[Project goal and research question.]


% =========================================================
\section{Dataset}

% Explain:
% - Where did the dataset come from?
% - How many records/patients are included?
% - What types of variables are available?
% - What is the target variable?
% - What counts as a 30-day readmission?

[Describe the dataset.]

[Describe the target variable.]


% =========================================================
\section{Exploratory Data Analysis}

% Include:
% - Basic descriptive statistics
% - Distribution of the target variable
% - Missing data
% - Important patterns found in the data
% - Useful graphs

\subsection{Descriptive Statistics}

[Add important descriptive statistics here.]


\subsection{Data Visualizations}

% Example:
% \begin{figure}[ht]
%     \centering
%     \includegraphics[width=\linewidth]{figure_name.png}
%     \caption{Description of the figure.}
%     \label{fig:example}
% \end{figure}

[Add graphs and explain the important patterns they show.]


% =========================================================
\section{Data Cleaning and Preparation}

% Explain:
% - How were missing values handled?
% - Were any columns removed?
% - Were any variables grouped or changed?
% - How were categorical variables encoded?
% - Were numerical variables scaled?
% - What did the final dataset look like?

\subsection{Missing Data}

[Explain how missing values were handled.]


\subsection{Variable Preparation}

[Explain which variables were removed, grouped, encoded, or changed.]


\subsection{Final Dataset}

[Describe the dataset after cleaning and preparation.]


% =========================================================
\section{Methods}

% Explain:
% - What models did we test?
% - Why did we choose them?
% - How did we split the dataset?
% - What hyperparameters/settings were used?
% - What software/libraries were used?
% - How did we evaluate the models?

\subsection{Models}

% Add only the models that were actually used.

\textbf{Model 1:} [Model name and short explanation.]

\textbf{Model 2:} [Model name and short explanation.]

\textbf{Model 3:} [Model name and short explanation.]


\subsection{Training and Data Split}

[Explain how the training, validation, and test data were created.]


\subsection{Hyperparameters}

% Example table -- change it later based on the actual models.

\begin{table}[ht]
\centering
\caption{Model Hyperparameters}
\begin{tabular}{ll}
\toprule
Model & Main Hyperparameters \\
\midrule
Model 1 & [Add values] \\
Model 2 & [Add values] \\
Model 3 & [Add values] \\
\bottomrule
\end{tabular}
\end{table}

\subsection{Evaluation Metrics}

% Explain the metrics actually used.
% Possible examples:
% Accuracy
% Precision
% Recall
% F1 Score
% ROC-AUC
% PR-AUC

[Explain how model performance was evaluated.]


% =========================================================
\section{Results}

% Include:
% - Results from each model
% - Which model performed best?
% - Important graphs
% - Confusion matrix / ROC curve if used
% - Comparison with previous work if appropriate

\subsection{Model Performance}

\begin{table}[ht]
\centering
\caption{Model Performance}
\begin{tabular}{lcccc}
\toprule
Model & Precision & Recall & F1 & ROC-AUC \\
\midrule
Model 1 & [ ] & [ ] & [ ] & [ ] \\
Model 2 & [ ] & [ ] & [ ] & [ ] \\
Model 3 & [ ] & [ ] & [ ] & [ ] \\
\bottomrule
\end{tabular}
\end{table}


\subsection{Main Findings}

[Which model performed best?]

[What did we learn from the results?]

[Were there any important predictors or patterns?]


% =========================================================
\section{Discussion}

% Explain:
% - What do the results mean?
% - Why might one model have performed better?
% - What did we learn about readmission?
% - What worked well?
% - What did not work?
% - What limitations did we find?

[Discuss the meaning of the results.]

[Discuss what worked and what did not work.]

[Discuss limitations of the project.]


% =========================================================
\section{Future Work}

% Explain:
% - What could be improved?
% - What other models could be tested?
% - What other data could be useful?
% - What questions remain unanswered?

[Future work idea 1.]

[Future work idea 2.]

[Future work idea 3.]


% =========================================================
\section{Conclusion}

% Briefly summarize:
% - What did we try to do?
% - What was the main result?
% - What was the main lesson from the project?

[Write the conclusion after the results are complete.]


% =========================================================
\section*{Author Contributions}

Siddique Chand: [Add contribution.]

Aubin Mugisha: [Add contribution.]

Cameron Letendre: [Add contribution.]


% =========================================================
\begin{thebibliography}{9}

% Add references here as the project develops.

\bibitem{uci}
UCI Machine Learning Repository,
``Diabetes 130-US Hospitals for Years 1999--2008.''

\end{thebibliography}

\end{document}